\documentclass[11pt]{article}
\usepackage[margin=1in]{geometry}
\usepackage{amsmath,amssymb,array,longtable}
\usepackage[hidelinks]{hyperref}
\newcommand{\Lfac}{\mathcal L}
\newcommand{\C}{\mathbb C}
\newcommand{\Q}{\mathbb Q}
\title{Factorial, Image and Rigidity: evidence, transfers and remaining work}
\author{Research audit}
\date{17 November 2026}
\begin{document}
\maketitle

\noindent\textbf{Current result.}
The complete critical-value argument for Furter rigidity has been checked
and written out in \href{RIGIDITY_CRITICAL_VALUES.tex}{the companion proof}.
It is credited to liqsweep's September 2026 manuscript, not presented as a
new discovery. It proves rigidity and SFC for the special family
$(\prod y_i)\sum\mu_i y_i$. Long's full binary FC proof and its pinned EIT I
dependency also pass the \href{LONG_EIT_AND_FC2_AUDIT.tex}{complete internal
written audit}. General FC in dimension at least three and general SFC
remain unproved here.
Unrestricted Image is already false, including the two-pair special case
retained in this repository. All new assurance flags for external review,
formal verification and independent computational replay remain false.

The \href{SIMPLEX_PENCILS_AND_AFFINE_LIMITS.tex}{all-dimensional simplex-pencil
proof} now also closes every phase $p(\ell(T))$ for one complex affine
coordinate $\ell$, with positive integer Dirichlet shapes. Its homogeneous
Gamma/factorial lift is $H(S,L)$, where $S$ is the radial sum and $L$ is
one arbitrary linear form. It retains all polynomial markers and repeated
node values; exact chain export and pencil recognition are implemented.
The same note proves algebraic-affine exponential nonvanishing and a flat
next-layer factorial class with arbitrary complex lower terms.

The \href{TWO_ANGULAR_COORDINATES.tex}{two-angular cubic proof} goes beyond
that pencil class. An explicit rank-two connection, elementary Riccati
obstruction, strict interior minimum and ordinary-point arithmetic lifting
give actual linear comparison on a triangle. Reduction at $h=0$ modulo five
then handles every triple of integer Dirichlet shapes at once. The resulting
degree-four Gamma/factorial families have eventually nonzero moments in
every dimension at the stated algebraic scales, with arbitrary complex
lower terms. The \href{two_angular.py}{exact Stokes and arithmetic exporter}
retains the actual boundary; a zero remainder makes no zero-period decision.

The \href{SFC_WINDOWS_AND_ROCKS_AUDIT.tex}{finite-window and Rocks audit}
now checks genuine SFC cases and implements all-index consecutive-window
certificates for linear powers, including every polynomial marker and
integer shape sums at colliding active nodes. The source identification,
common-prime hypothesis and affine draft window are corrected below.

\section*{Statements that must be distinguished}
FC$(r)$ says that $\Lfac(f^k)=0$ for every $k\geq1$ forces $f=0$, where
$\Lfac(y^a)=\prod_i a_i!$. SFC asks that every block of $\mathcal N(f)$
consecutive positive exponents contain a nonzero moment, for every
nonzero $f$. This is stronger than absence of an infinite zero tail.

The special Image space is
\[
 M_r=\sum_{i=1}^r(\partial_{z_i}-\xi_i)\C[\xi,z]=\ker\mathcal E_r,
 \qquad \mathcal E_r(\xi^a z^b)=\partial_z^a z^b.
\]
The claim that this is Mathieu--Zhao is false for $r=2$, as the retained
proof below shows. The broader regular-sequence Image conjecture includes
this instance (take $A=\C[\xi]$, $a_i=\xi_i$), so it is false as well.
Image for unrelated locally nilpotent derivations is a different assertion;
one must specify the operator and coefficient algebra.

\section*{Established results and current claims}
The original source
\href{https://arxiv.org/abs/1008.3962v2}{van den Essen--Wright--Zhao (2011),
Section 4} proves FC for one variable, two monomials, a monomial times a
polynomial with nonzero constant term, powers of linear forms and diagonal
equal-power sums. Two monomials need only the first two moments.
Its Proposition 4.10 gives a further coordinate-divisibility support class.
Real polynomials up to a nonzero complex scalar satisfy FC immediately
by positivity of the second moment; the difficulty is complex cancellation.

\href{https://doi.org/10.1017/S0004972719000546}{Liu--Sun (2020)} proves FC
for every homogeneous binary polynomial. This covers uniform interval
integration after homogenization, not arbitrary inhomogeneous binary FC,
and not triangle integration.

\href{https://arxiv.org/abs/1304.3956}{Edo--van den Essen (2014)} introduces
SFC and proves its exact connection with rigidity on the special family.
It records published $R(1)$ and $R(2)$ and proves ordinary FC for that
family in every dimension. Its two-monomial analysis is not a proof of
unrestricted SFC. The linked
\href{https://people.math.sc.edu/rocks/BradyRocksSFC2015.pdf}{Rocks (2015) PDF}
is a 28-page author article preprint, not his separately catalogued
\href{https://doi.org/10.7936/K7BR8QB4}{dissertation}.
The journal article is \href{https://doi.org/10.1016/j.jalgebra.2017.01.039}{Rocks,
\emph{On the incompatibility of diophantine equations arising from the Strong
Factorial Conjecture}, J. Algebra 478 (2017), 341--352}.
The \href{SFC_WINDOWS_AND_ROCKS_AUDIT.tex}{core draft proof audit} now checks
powers of linear forms, powers of sums of powers of one COMMON prime,
and adjunction of new disjoint linear variables. The affine proof gives
the correct $m+1$-moment window; its displayed draft statement says $m$
and is refuted by $Y-1$. The draft's constant-added prime-power result is
limited to starts $n+m<p$. These corrections are not attributed to the
unavailable final publisher text. Other draft two-monomial claims and
the full dissertation remain unaudited.

\href{https://doi.org/10.4064/cm6761-3-2016}{Filaseta--Rocks (2016),
Colloq. Math. 145, 307--314} proves a fixed-degree lower density at least
$\log2$ of irreducible moment polynomials $f_{n,d}$. The exact identity
$n f_{n-1,d}=n f_{n,d}-x f'_{n,d}$ identifies adjacent common roots with
multiple roots. The density statement is useful evidence for two-term
SFC, not an all-index result or a general solution.

The \href{PRIME_POWER_TWO_TERM_SFC.tex}{prime-power two-term proof} now
closes $A+BY^{p^a}$ at EVERY starting index, for every prime $p$ and
$a\ge1$. Scaled Newton edge residuals are reversals of truncated
exponentials of degree below $p$, so the parameter moment polynomials
are squarefree; the adjacent derivative identity excludes two consecutive
zero moments. Integer Gamma shapes $1\le\alpha\le p^a$, a specified
low-degree marker class and independent-linear adjunction are included.
The ordinary pair ideal has length $n(n+1)$ and sharp coefficient
nilpotence exponent $2n$ over every $\Q$-algebra. This is an internal
written theorem with no priority claim; the first non-prime-power exponent
left by this argument is six. General two-term and unrestricted SFC
remain unresolved here.

The \href{MIXED_PRIME_RESIDUAL_WINDOWS.tex}{mixed-prime residual proof}
now gives an exact insertion rule for $d=mp^a$ and unbounded indexed
windows at terminal $N=kp^r$, under $p>mk$ and a squarefree small kernel.
It also combines possible square-divisor degrees across primes, certifying
some exponent-six windows where neither prime alone suffices. All-starting-
index exponent-six SFC remains open; a successful finite table or indexed
family is not substituted for it. The journal Online First version of
Filaseta--Rocks was located and its needed single-edge lemma checked.

The \href{HYPERGEOMETRIC_MONOMIAL_WINDOWS.tex}{uniform monomial-window proof}
now handles every positive real Gamma/Dirichlet shape and every polynomial
marker for nonconstant $a+bZ$, with $Z$ a monomial of total degree $D$.
An exact polynomial $R(t)$ decides a whole marked zero sequence or gives
a $D+\deg R+1$ window at every start. The extra simplex singularity is
checked by an indicial exponent exceeding the moment polynomial's degree.
Unmarked windows have length $D+1$; actual expanded block-sum support
therefore proves SFC whenever it covers that bound, for example
$a+b(Y_1+Y_2)^6$ with eight monomials. Ordinary coefficient ideals and
original Gaussian markers are retained. For the two-monomial $a+bY^6$,
seven does not replace the required two; its full SFC remains open.

The \href{COUPLED_BLOCK_PERIODS.tex}{coupled-block proof} now checks the
complete EIT VI block-radial limit and extends it to every positive real
Gamma shape and every polynomial marker. A genuinely coupled angular
product on $r$ intervals has period $J_{r-1}(\xi)/\xi$, nonzero at every
algebraic scale by an explicit independent $E$-function system. Exact
algebraic-marker period tests retain zero leading limits separately.
The resulting nonflat families in $2r$ factorial variables have eventual
nonzero moments at algebraic scales, and FC with infinitely many surviving
moments at every complex scale. Ordinary universal tail ideals are unit
over every $\Q$-algebra. Arbitrary coupled triangle phases and unrestricted
FC/SFC remain open; the remaining geometric/elliptic source inputs are
not covered by this block-limit audit.
Its singleton-block corollary covers every single nonconstant leading
monomial, allowing zero exponents and arbitrary complex lower layers in
additional variables. In particular,
$F=X^3+Q_{\leq2}(X,Y,Z)$ has
$\mathcal L(F^m)/(3m)!\to\exp(Q_2(3,0,0)/27)\ne0$.
The same inactive-variable extension preserves every proved nonzero
block limit when $G(X,0)$ is in the original family. Markers use active
degree and unchanged inactive Gamma moments. No algebraic-scale
restriction is needed for the single-top-monomial conclusion.

\href{https://doi.org/10.1112/jlms/jdu064}{Furter (2015)} relates composition
rigidity to plane automorphism polydegree closures. The correct citation is
\emph{J. London Math. Soc.} 91(1), 180--202, DOI jdu064. Some recent
secondary records instead give 91(3), 743--760; do not propagate that error.
\href{https://arminstraub.com/downloads/pub/polydegree-conjecture.pdf}{Lewis--Perry--Straub (2019)}
gives algorithmic sufficient conditions and finite computations.
\href{https://github.com/ipitchford/furter-r3-through-299}{A 2026 unrefereed
finite certificate package} claims $R(3)_n$ through $n=299$; it is not a
universal argument and was not replayed here.

\href{https://github.com/blueberryvertigo/polynomial-composition-rigidity/tree/051c74aad9ef9bf8804c371964803aee099c5b63}{Liqsweep (September 2026)}
claims full rigidity via attracting fixed points and critical values.
The companion note checks the whole analytic proof, proves the inverse
index conversion, and transfers it to the special factorial/simplex family.
Thus the internal mathematical result is now full rigidity; peer review
and independent formalization are not claimed. We do not independently
audit the manuscript's scheme-theoretic composition-map or polydegree
corollaries in this turn.

\section*{Long's two-variable manuscript: actual dependencies}
The current \href{https://github.com/long-mathematics/factorial-conjecture-two-variables/tree/8b1eb46ff530122ed25f588b78bf9fdc98cb9577}{pinned manuscript}
claims full inhomogeneous FC$(2)$. Its nonflat leading-part argument uses
Mittag--Leffler continuation, polynomial monodromy, block traces and
Lindemann--Weierstrass. The flat case additionally invokes
\href{https://github.com/long-mathematics/exponential-integral-theorem/blob/22d98c4f27e606b5b25fa673c33f9c56c53cc86b/papers/01-exponential-integral-theorem/exponential-integral-theorem.tex}{EIT I, pinned Corollary 1.4}:
$\int_a^b e^{\xi q(t)}dt\ne0$ for distinct algebraic endpoints, nonconstant
algebraic polynomial $q$, and nonzero algebraic $\xi$.
Both manuscripts are unrefereed. Their main proof chains have now been
checked in \href{LONG_EIT_AND_FC2_AUDIT.tex}{the detailed audit}, including
EIT's semisimplicity, rational-projector constancy, Beukers lifting, and
tree-density argument, and FC's corrected leaf transport, rational path
module, simultaneous block traces and positive conjugate pairing.
The exact published inputs have been matched to their uses; no generic
critical-value assumption or absolute irreducibility is inserted.
The internal registry records this mathematical verification with all
external-review and formal-verification flags false. The 355 new exact
controls check transcription, not the analytic universal quantifiers.

Here is the flat limit, independently checked algebraically. Write
\[
 f(rt,r(1-t))=r^D+\sum_{j=1}^D A_j(t)r^{D-j}.
\]
In the expansion of the $n$th power, put $K=\sum_j k_j$ and $J=\sum_jjk_j$.
The normalized radial coefficient is
\[
 C_{n,k}=\frac{n!}{(n-K)!}\frac{\Gamma(Dn-J+2)}{\Gamma(Dn+2)}
 \quad(K\leq n),
\]
and zero otherwise. Since $K\leq J\leq DK$, the first $K$ factors of the
denominator dominate the numerator, so $0\leq C_{n,k}\leq1$.
For fixed $k$, $C_{n,k}\sim D^{-J}n^{K-J}$.
The bound $\prod_j\|A_j\|^{k_j}/k_j!$ is summable over $k$, so dominated
convergence on $[0,1]$ gives
\[
 \frac{\Lfac(f^n)}{\Gamma(Dn+2)}\longrightarrow
       \int_0^1e^{A_1(t)/D}\,dt.                            \tag{3}
\]
If all pure moments vanish, the limit must vanish. For algebraic $f$ and
affine $A_1/D=\alpha t+\beta$, the right side is $e^\beta$ if $\alpha=0$,
and $e^\beta(e^\alpha-1)/\alpha$ otherwise, nonzero by Hermite--Lindemann.
For complex coefficients one first specializes the all-moment polynomial
equations to an algebraic point, retaining the flat normalization and the
affine condition by rational equations. Noetherianity and Nullstellensatz
justify this specialization; a finite sampled prefix alone would not.
This verifies the flat affine-next-layer subcase in every total degree,
including the flat degree-at-most-two case, without EIT.
The completed EIT audit handles an arbitrary algebraic next layer and
closes the general flat binary case. Algebraic specialization then gives
FC$(2)$ for complex coefficients. This does not assert a nonzero radial
limit for every transcendental coefficient choice: for
$f=(x+y)^2+4\pi i x$ that limit is zero although $\Lfac(f)\ne0$.

\paragraph{Additional higher-dimensional source.}
Long--Le Blanc's \emph{Exponential Periods and Factorial Moments}, EIT VI,
\href{https://github.com/long-mathematics/exponential-integral-theorem/blob/92db79eedeb064694ff20e9201d73fc2bc0857f4/papers/06-exponential-factorial-moments/exponential-factorial-moments.tex}{October 2026 source},
claims all-dimensional radial limits, explicit flat families with eventually
nonzero moments, degree-four families at a fixed algebraic complex scale,
and block-radial nonflat families. Its nonvanishing input comes from
Paper V and has additional hypotheses. Its finite gcd range is not a
proof at every parameter, and its complex-scale statement distinguishes
infinitely many nonzero moments from an eventually nonzero tail.
Its radial proof and the needed Paper V linear/extremum mechanism have now
been checked. The cubic two-coordinate case is proved independently in
\href{TWO_ANGULAR_COORDINATES.tex}{the companion note}, with a direct Riccati
irreducibility argument. The general geometric/Galois, elliptic
continued-density and nonflat block-family inputs remain outside this
completed audit. The source claims neither unrestricted FC nor arbitrary
simplex exponential nonvanishing.

\section*{Image: reuse the already stronger retained counterexample}
The registry already contains the two-pair witness in
\href{../../foundations/extended-geometry/TWO_PAIR_IMAGE_MATHIEU_COUNTEREXAMPLE.md}{the canonical retained proof}.
Set
\[
 R=\xi_1z_1+\xi_2z_2,\quad Z=\xi_1z_2,\quad
 W=2\xi_2z_1,\quad T=\xi_1z_1-\xi_2z_2,
\]
and $F=(R+Z)(R^2W-(2R+Z)T^2/2)$, $Q=Z$. Then for every $m\geq1$,
\[
 \mathcal E_2(F^m)=0,\qquad
 \mathcal E_2(QF^m)=\frac{(4m+2)!m!}{(2m+1)!!}\ne0.          \tag{4}
\]
The all-order proof uses circular Gaussian radial/angular separation and
$J_m'(X)=(1-X^2)^m$: the Taylor coefficients through degree $m$ of
$X^{m-1}J_m(X)$ at $X=1$ coincide with those of $J_m(1)X^{m-1}$.
Thus this is an all-order algebraic obstruction, not an inference from
finitely many moments. SIC$(1)$ holds, so two is the minimum failing
pair count. This witness is not in the independent radial invariant
subalgebra $\C[\xi_1z_1,\xi_2z_2]$. Diagonal extraction does not commute
with taking powers, and does not disprove factorial FC.
The five-pair Dvorsky--Long witness is a simpler separately retained
counterexample, but its dimension is not the sharp current endpoint.

\section*{What the Gaussian machinery transfers exactly}
For $Y_i=\frac12\sum_{j=1}^{k_i}X_{ij}^2$, put
$\alpha_i=k_i/2$ and $A=\sum_i\alpha_i$. Then
$\mathbb E Y^b=\prod_i(\alpha_i)_{b_i}$. With $R=\sum_iY_i$ and $T=Y/R$,
$T\sim\operatorname{Dir}(\alpha)$ is independent of $R$.
For lifts $F=R^Dp(Y/R)$ and $H=R^eq(Y/R)$,
\[
 \mathbb E[HF^m]=(A)_{e+Dm}\mathbb E[q(T)p(T)^m].             \tag{5}
\]
Orthogonal averaging over $\prod_i O(k_i)$ retains every Gaussian
polynomial marker and turns it into a polynomial in $Y$; splitting its
homogeneous pieces proves both directions of eventual marked vanishing.
The complete prior proof is
\href{../cross-programme-connections/GAUSSIAN_GAMMA_DIRICHLET_AND_OUTWARD_TRANSFERS_2026_11_16.tex}{the Gamma--Dirichlet transfer note}.

The new \texttt{moment\_bridges.py} implements exact moment functionals,
homogenization, Gaussian lifts, translation of arbitrary Gaussian monomial
markers by their explicit Reynolds factors, and input serialization for
the existing \texttt{gaussian\_source\_hom\_producer.py}. Its finite checks
call that producer's actual input parser. No expensive new CAS campaign
or universal pure-null decision is claimed.

Full binary FC gives a further exact structural consequence: in the entire
Gaussian subalgebra
$\C[(X_1^2+X_2^2)/2,(X_3^2+X_4^2)/2]$, eventual pure vanishing and
all-marker safety both force the phase to be zero. Over any commutative
$\Q$-algebra they instead characterize coefficientwise nilpotence.
For a binary factorial phase with coefficient ideal $J$, the scalars
making every actual marked period polynomial form the ordinary ideal
$\Gamma_J(B)=\bigcup_k\operatorname{Ann}_B(J^k)$.
The proof in \href{LONG_EIT_AND_FC2_AUDIT.tex}{the completed audit} retains
ordinary nilpotents and works without a Noetherian hypothesis.

Safety (eventual marked vanishing for every fixed marker), purity
(every positive pure moment vanishes) and a short consecutive zero block
are distinct requirements. The current whole-source Gaussian producer
decides safety, and tests purity inside its safe output. An unsafe output
alone does not refute purity. A separate valid unit-period annihilator
with actual initial jets can certify a pointwise pure decision; building
that dedicated adapter remains useful work. General Gaussian MZ is false,
so purity cannot be upgraded to safety without a theorem for the radial
subalgebra. Neither the Gaussian counterexample nor the Image witness
automatically belongs to that subalgebra.

\section*{All ordinary constant-plus-power two-term families are completed}
The \href{ALL_UNIVARIATE_TWO_TERM_SFC.tex}{uniform effective simplicity proof}
now establishes
\[
 \boxed{F_{n,m}(t)=\sum_{k=0}^n\binom nk(mk)!t^k
        \text{ has simple complex roots for all }n,m\ge1.}
\]
Hence every nonzero $A+BY^m$, with arbitrary complex coefficients and
positive integer $m$, satisfies Strong Factorial at every starting index
in ordinary Gamma shape one. Reversal preserves multiplicities and
$P'_{n+1,m}=(n+1)P_{n,m}$ gives the exact consecutive-moment equivalence.
This covers one variable with a constant term; it is not unrestricted SFC,
other Gamma shapes, general monomial pairs, or general Factorial.

For $n\ge10000$ and $6\le m\le\sqrt n$, put $A_m=(m/e)^m$ and
$v=z/(n^mA_m)$. Every root lies in $1/2<|v|<21/10$.
The scaled integral splits exactly at $s=m/2$. On the full disk
$|v|\le11/5$, its tail has a single-saddle expansion with analytic
relative error $<22/\sqrt{mn}$, complete contour/tail control and an
effective Cauchy derivative estimate. Its endpoint has a uniform Watson
estimate on $|v|\ge2/5$. At every possible cancellation, the logarithmic
derivative gap is $>n/3-4>0$, excluding every multiple root.
The large-exponent coefficient criterion covers $2^m\ge4n^2$; otherwise
$m<2\log_2(2n)<\sqrt n$. Prime powers include all smaller exponents.
Thus $10000$ is an effective simplicity cutoff uniform over all $m$,
with the root annulus asserted only in the stated analytic range.

Below the cutoff, the retained prime-power and
\href{EFFECTIVE_EXPONENT_SIX_SIMPLICITY.tex}{all-index exponent-six proofs}
are reused. All $92935$ further necessary pairs have
\href{all_exponent_prefix_certificates.json}{explicit modular B\'ezout certificates}:
$92900$ preceding-prime kernels have degree at most $35$, and the other
$35$ full-degree reductions have index at most $126$. Every prime,
identity and pair-domain coverage is checked exactly. The real-place
parity and shifted-Newton lemmas complete the finite proof.
The \href{uniform_two_term_sfc_verification.json}{uniform verifier}
checks $92990$ controls, including every additional finite pair; the
universal evidence is the analytic proof. Independent, formal and external
assurance remain false; no numerical root scan or priority claim is made.

For this sequence the ordinary exceptional sets are now \emph{empty}:
$E_m(X)=D_{\rm all}(X)=0$. The earlier density bounds and MR source audits
are retained as historical proof routes, but the new pointwise proof uses
no analytic-number-theory input. The prior
\href{EFFECTIVE_SADDLE_ANNULUS.tex}{saddle-annulus result} is retained as a
completed step, with its annulus-only $10^6$ threshold distinguished from
the actual uniform simplicity cutoff $10000$.

\section*{The broader scope and the propagation boundary}
The overarching goal remains a proof or counterexample for general
Factorial and/or unrestricted Strong Factorial. The former
constant-plus-power all-index gate is closed directly. The next support
family must be genuinely outside that closure; for instance
$f=x^a(A+Bx^d)$ has moments
\[
 \sum_{k=0}^n\binom nk A^{n-k}B^k(an+dk)!,
\]
with Gamma shape $an+1$ varying with $n$. The fixed-shape Appell derivative
identity is unavailable, so the current squarefreeness equivalence does
not automatically prove its consecutive-moment assertion. Retained
literature must be checked before treating any such family as a new open gate.
Genuinely multivariate monomials and arbitrary polynomial supports also
require their own arguments; no reduction of unrestricted FC/SFC to the
completed one-variable family is asserted.

Before the pointwise closure, one possible use of the density bounds
would have required a theorem forcing enough distinct bad terminal indices
\emph{inside the counted family}. Such a theorem was not established.
Now even one valid transfer of a general SFC failure to a genuine bad
$F_{n,m}$ would contradict the pointwise theorem, but no such transfer is
supplied by the retained comparison.
The existing
\href{../moment-vanishing-audit/POLYNOMIAL_RECOMPOSITION_REFLECTION_AND_UNIVERSAL_FAILURE.md}{failure-preserving recomposition theorem}
assumes an all-order pure unsafe Gaussian or intrinsic Haar seed, or its
stated eventual-purity variant. A finite SFC zero block does not provide
those hypotheses. The safety-reflection theorem supplies no missing
finite-block conversion either, and recomposition may leave the counted
family. These hypotheses remain distinct. No new general propagation
or arbitrary-support reduction is claimed.

\subsection*{All adjacent monomial supports are completed}
The \href{UNIFORM_PADE_PARABOLA_SFC.tex}{uniform shared-parabola proof}
now establishes Strong Factorial for every $X^a(A+BX)$, every integer
$a\ge0$, every starting index, and every positive Gamma shape $\beta$.
It also covers $\prod_j X_j^{a_j}(A+BX_1)$ under independent positive
Gamma shapes, because the other variables contribute nonzero factors.
The pinned Rocks 2015 author manuscript gives earlier partial ranges;
no priority or complete final-publication audit is claimed.

Define $R_N(t,z)=\sum_{k=0}^N(-1)^k\binom Nk(t+1)_kz^{N-k}$.
Its complex symmetric tridiagonal matrix places every zero, for every
real $t>-1$, strictly inside
\[
                   (\operatorname{Im}z)^2<4N(\operatorname{Re}z+N).
\]
A common zero of $R_N(l+a,z)$ and $R_{N-1}(l,z)$, where $l>-1$ and
$a$ is a positive integer, would also be a zero of an explicit
parameter-shift continuant. The normalized recurrence has
$b_j(1-b_{j-1}/c_j)\ge N$. An elementary continued-fraction invariant,
proved by completing a square, makes that continuant nonzero in the
same parabola. These complementary inclusions prove coprimality
uniformly, with no cutoff or finite prefix.

The identity
$R_n(\beta+an-1,z)=z^n\sum_k\binom nk(\beta+an)_k(-1/z)^k$
transfers coprimality to the two consecutive moments.
The \href{pade_parabola_sfc_verification.json}{focused receipt} checks
464 exact controls. This algebraic proof is independent of the earlier
contour, density, irreducibility, parity and prime-gap stages.
The earlier \href{PADE_ADJACENT_MONOMIALS.tex}{finite-$a$ strip proof}
and its 471 controls remain preserved evidence. Its stronger uniform
vertical spectral bound is still unproved, but is unnecessary for the
now completed adjacent-support assertion.

The next support gate must have a genuine gap $d\ge2$ in
$X^a(A+BX^d)$ with $a\ge1$, or more general multivariate supports.
No reduction of unrestricted FC/SFC to the completed families is
established. All independent, formal and external assurance flags
remain false; the overarching conjectures remain open.

\subsection*{A completed nonadjacent dyadic family}
The \href{DYADIC_BINOMIAL_SFC.tex}{dyadic unit-transfer proof} establishes
\[
 X^q(A+BX^q)\text{ satisfies SFC at every start},\qquad
 q=2^r,\ r\ge0,\quad1\le\beta\le2q\text{ an integer}.
\]
For $q\ge2$ the gap is genuinely nonadjacent. The normalization uses one
variable scaling $2^{-q}$ at both orders, and odd factorial units. The
resulting integral moment polynomials retain degree and reduce to Bessel
$Y_n$ or $D_n=(n+1)Y_n+zY_n'$. Both satisfy the Fibonacci recurrence
modulo two, so every consecutive resultant is odd. Arbitrary rational
unit sequences multiplying $(L+h)!$, $h=0,1$, admit the same transfer.
No prior contour, density, parity or prime-gap result is needed.

For the raw Gamma moment polynomial
$M_n(t)=\sum_k\binom nk(\beta)_{q(n+k)}t^k$, the proof also gives
\[
 v_2\operatorname{Res}(M_n,M_{n-1})
 =3qn(n-1)-(n-1)s_2(n)-ns_2(n-1),
\]
where $s_2$ is binary digit sum. This is exact at every index, independent
of the particular integer shape in the stated range.
The \href{dyadic_binomial_sfc_verification.json}{focused receipt} records
2784 exact controls; the proof, rather than finite replays, is universal.
Other-variable monomial factors are allowed, with any positive independent
Gamma shapes.

The former first remaining one-variable gate was $X(A+BX^2)$ in ordinary shape
one; it is now closed below, at all positive shapes. More generally, for every gap $e\ge1$, the new coefficient identities
$A_{n,e}=T_{n,e}+(e/n)\theta T_{n,e}$ and
$A_{n-1,e}=T_{n,e}-(1/n)\theta T_{n,e}$ identify its two-moment
obstruction EXACTLY with a multiple root of
$T_{n,e}=\sum_k\binom nk(n)_{ek}t^k$. This is an equivalence;
all-index simplicity is not asserted.

For the cubic and every prime $n<q\le3(n-1)/2$, $n\ge6$, put $d=q-n$.
Primitive positive-rise Newton edges confine every rational square divisor
$g^2\mid T_{n,2}$ to even degree $\deg g\le2\lfloor d/8\rfloor$.
It reduces degree-preservingly into a square divisor of the explicit
horizontal kernel $K_d$. Equivalently, the raw moment gcd reduces into
$\gcd(\overline H_d,\overline H_{d+1})$, with
$K_d=(H_d+2H_{d+1})/3$ and
$\theta K_d=d(H_{d+1}-H_d)/3$.
In particular an admissible prime at distance at most seven certifies
the pair. Uniform prime/kernel coverage remains open; the prime-window
lemma does not assert prime existence at every small index.
 At terminal index three its two Newton polygons share slope $3/2$
and residual $1+Y$, although the rational polynomials are coprime.
Similarly the $h=2$ shifted Bessel reductions at indices two and three
both equal $z^2$. These are limitations of the residue certificate,
not SFC counterexamples or an extension to every shape. Exact examples
are retained in the \href{dyadic_binomial_frontier_diagnostics.json}{boundary record}.
General gaps, arbitrary supports, general FC and unrestricted SFC remain
open; all assurance flags remain false.

\subsection*{Tilted cubic prime-power indices and exact fold geometry}
The \href{TILTED_CUBIC_PRIME_POWER_AND_FOLD.tex}{prime-power/fold proof}
now proves $T_{p^r,2}$ irreducible over $\mathbb Q_p$ and $\mathbb Q$
for EVERY prime $p$ and $r\ge1$. These are prime-power INDICES, distinct
from the earlier prime-power exponent theorem. For $p\ge5$ the polygon
is a single edge with denominator $p^r$; separate binary and ternary
arguments give primitive heights $2p^r-1$ and $p^r-1$.

More generally, $n=ap^r+i$, $p>3\max(a,i)$, has horizontal width $i$
and slope $2(p^r-1)/(p^r(p-1))$. The horizontal residual is $T_i$;
the sloped residual has Gamma shape $a$ at offset zero and shape $a+1$
at positive offsets. Parity and a nonsquare leading coefficient close
ten complete shifted index families listed in the proof. This remains
indexed coverage from this arithmetic method; the subsequent adjoint proof
below closes the whole cubic at all positive Gamma shapes.
The \href{tilted_cubic_frontier_verification.json}{1294 exact controls}
check these formulas without extrapolating a finite range.

For the analytic alternative, the exact saddle cubic has folds
$w=1/6\pm i\sqrt3/18$, with nonzero third and mixed derivatives.
No contributing-contour theorem, Airy estimate, root-location consequence
or effective simplicity cutoff is inferred from these data.
The retained \href{tilted_cubic_frontier_diagnostics.json}{exact diagnostics}
show that fixed primes $2,3,5,7,11$ leave degree six at index 732,
while thirteen removes it, and two simple large-prime exclusion guesses
fail at indices 51 and 118. Also $K_{11}$ has degree five but only three
real roots, all positive; a blanket real-rooted-kernel shortcut is false.
None is a repeated-factor counterexample. Uniform arithmetic coverage and
a contour-based simplicity theorem remain unproved as method statements;
the whole cubic SFC gate is closed independently below.

\subsection*{Whole cubic closure by explicit adjoint separation}
The \href{CUBIC_ADJOINT_STABILITY_SFC.tex}{new adjoint proof} proves
\[
 T_{n,\gamma}(t)=\sum_k\binom nk(\gamma)_{2k}t^k
 \quad\hbox{squarefree for every }n\ge1,\ \gamma>0,\ \gamma\ge n-2.
\]
After quadratic reversal the monic third-order equation has explicit
adjoint solutions $e^{\mp z}z^{1-\gamma-2n}B(\pm z)$, where
$B(z)=\mathbb E(z+U+V/2)^n$ for independent Gamma shapes
$\gamma,n+1$. The rederived Gamma tridiagonal bound and classical
Grace--Walsh--Szeg\H{o} half-plane polarization show $B$ strictly Hurwitz
in the stated range. Both Lagrange concomitants vanish, by a rational
logarithmic derivative identity. A repeated root would force
$B(z)=B(-z)=0$, contradicting their half-plane separation.

With $\gamma=n+\beta-1$, exact two-by-two coefficient identities now
close $X(A+BX^2)$ SFC at EVERY start for EVERY positive Gamma shape
$\beta$. Other-variable monomial factors are allowed. The
\href{cubic_adjoint_stability_verification.json}{271 exact controls}
check the algebra and regression cases. The universal proof has no
cutoff, finite prefix, parity or prime-coverage requirement. Its only
imported theorem is the explicitly sourced classical polarization input.
All independent/formal/external assurance and novelty flags remain false.
Other lower exponents/gaps, general supports and unrestricted FC/SFC
remain open.

\subsection*{All-gap adjoint structure and the precise rank frontier}
The \href{HIGHER_GAP_ADJOINT_RANK.tex}{higher-gap proof} constructs
$d$ independent adjoint solutions $e^{-\omega z}z^{1-\gamma-dn}B(\omega z)$,
$\omega^d=1$, for EVERY gap $d\ge2$. The auxiliary satisfies
$B_{n+1}=(z+dn+\gamma+d-1)B_n-zB_n'$.
A repeated root is exactly a rank defect of the explicit $d$ by $(d-1)$
polynomial adjoint jet matrix. A full Wronskian determinant reconstructs
the reversed moment polynomial. These are all-parameter identities,
not a uniform rank-exclusion theorem.

For gap three, a common rotation orbit of a real pair-Wronskian $W$
is the obstruction. Uniform coprimality or strict interlacing of its
three cyclic sections would suffice, as would zero-freeness in the
CLOSED sector $|\arg z|\le\pi/3$.
The stronger full Hurwitz assertion is false already at $n=\gamma=4$:
the seventh determinant is strictly negative. The
\href{higher_gap_adjoint_rank_verification.json}{138 exact controls}
include actual multiple roots outside positive Gamma scope, and exact
section interlacing ONLY at $n=\gamma=4,6,8$. No uniform interlacing,
sector or all-index higher-gap SFC conclusion is inferred.
The same note proves the uniform initial block $n\le4$ at EVERY gap
and EVERY positive Gamma shape. A quartic moment square would force
zero third central moment, whereas increasing convexity of $(x)_d$
gives the strict lower bound $\kappa_3\ge2a(b-a)^2>0$.
The extension imports no external theorem. General FC/unrestricted SFC
remain open, and all independent/formal/external assurance flags are false.

\subsection*{A rational Wronskian recurrence and a weaker real-section gate}
The \href{GAP_THREE_WRONSKIAN_RECURRENCE.tex}{new gap-three proof}
derives $B_{n,\gamma-1}=B_{n,\gamma}-B_{n,\gamma}'$ and identifies $W$
with a divided difference of consecutive auxiliaries. With
$a=3n+\gamma+2$ the exact recurrence is
\[
 W_{n+1,\gamma}=a(a-1+z)W_{n,\gamma}-azW_{n,\gamma}'
                                         +z^2W_{n,\gamma-1}.
\]
Centering $\gamma=n+b-1$ gives a rational triangular exporter and
nonnegative coefficients in BOTH $z$ and $b$.

More significantly, the three cyclic sections have NO common real root
at any positive Gamma shape, by the moment integral and the exact rank
bridge $t=-1/u$. Thus ONE nonzero real-rooted section or real constant
linear combination suffices for simplicity. Its choice can depend on
$n,\gamma$. Strict interlacing of all sections is unnecessary.
At $n=5$, a quadratic section has a discriminant with eleven strictly
positive coefficients in $\gamma$, closing the initial block through
terminal five at gap three for EVERY positive shape.

The \href{gap_three_wronskian_verification.json}{145 exact controls}
include a formal recurrence step whose separately interlacing,
positive, monic and strictly log-concave inputs produce a nonreal
section. Those inputs are not the actual common-Gamma pair. This
refutes that full-interlacing induction shortcut, not SFC or the weaker
one-section target. Uniform real-section existence and general
FC/unrestricted SFC remain open; all assurance flags remain false.

\subsection*{A uniform polar companion and a nonzero-pivot alternative}
The \href{GAP_THREE_POLAR_AND_CONTINUANT.tex}{new companion proof}
writes the auxiliary coefficients as backward differences of
$p_n(\gamma)=\prod_{j=0}^{n-1}(\gamma+2+3j)$.
The binomial companion $J$ and $K=\sum n_{\underline j}b_jz^j$ have
simple real roots strictly below $-1$ for ALL $\gamma>n-2$.
An elementary polar-derivative proof is independent of the related
finite-difference source. Its infinite-ray theorem cannot be applied
to backwards sampling, which has positive roots on its forward ray.
A generic inverse-companion example at degree eight has no individually
real-rooted Wronskian section, so the actual arithmetic-progression
Gamma structure remains essential to any transfer argument.

Alternatively normalize the three sections as $q_{2n-r}=W_r/W_r(0)$
and descend by
$\lambda_j=[u](q_j-q_{j-1})$,
$q_{j-3}=(q_j-q_{j-1})/(\lambda_ju)$.
If every pivot is nonzero, descent to $q_0=q_1=q_2=1$ proves gcd one.
No positivity is required. A zero pivot is inconclusive, as an explicit
coprime positive triple demonstrates. Uniform actual-family nonzero
pivots have not been proved. The
\href{gap_three_polar_continuant_verification.json}{84 exact checks}
retain formal parameter conditions, tamper rejection and finite controls.
At the native prime three, residuals already repeat at index22, and
residual occupancy plus parity still permits degree two at63. These are
method limitations, not SFC counterexamples or all-prime limitations.
General FC and unrestricted SFC remain open; all assurance flags are false.

\section*{Remaining programme, without reducing the goal}
\begin{enumerate}
\item Finish EIT V geometric/Galois and continued-density inputs and
      the remaining higher odd-degree and elliptic angular inputs in VI.
      The complete block-limit/product step and coupled-product extension
      are now checked; retain parameter, amplitude and finite-range boundaries.
\item Prove or refute general FC and SFC. General triangle integration is
      equivalent to homogeneous FC$(3)$; tetrahedra correspond to FC$(4)$.
      The special simplex family proved here is a strict subclass.
\item Implement original-unit period certificates that test purity even
      when all-marker safety fails, using original IBP and certified jets.
\item Find a structural obstruction or classification for pure-null
      phases invariant under all independent block rotations. The exact
      Gamma bridge reduces FC to this restricted question, without solving it.
\item Complete bibliography review (including Rocks and subsequent
      operator-specific Image papers). This is a sourced working survey,
      not a claim that every paper has been read or that priority is settled.
\end{enumerate}
The requested inactive-variable extension is proved in the coupled-block
note. The general FC/SFC programme remains open. Rigidity and the binary Factorial
claim now have complete internal proof audits, and unrestricted Image
has a retained counterexample. The general higher-dimensional Factorial
and unrestricted Strong Factorial questions still require a resolution.
\end{document}
